\documentclass[12pt]{article}
\usepackage[T1]{fontenc}
\usepackage[utf8]{inputenc}
\usepackage{lmodern}
\usepackage{textcomp}
\usepackage{enumitem}
\usepackage{geometry}
\geometry{margin=1in}
\begin{document}
\begin{center}
Answer Key - Constitutional Law - Undergraduate Level Assessment 16 \
\small (Answers are ordered exactly as in the question paper.)
\end{center}

\section*{Multiple Choice}
\begin{enumerate}[label=\arabic*.]
\item \textbf{Answer:} D
\\ \textit{Options:} A) They adopt different attitudes towards the role of the courts in maintaining order.; B) They disagree about the role of law in society.; C) They have opposing views about the nature of contractual obligations.; D) They differ in respect of their account of life before the social contract.
\\ \textit{Note:} Hobbes posits a chaotic state of nature characterized by constant fear and competition, necessitating a strong sovereign for order, while Locke envisions a more peaceful and rational state of nature where individuals possess natural rights, thus highlighting their differing views on life before the social contract.
\item \textbf{Answer:} B
\\ \textit{Options:} A) It can be remedied by redistribution of wealth.; B) If each person's holdings are just, then the total distribution of holdings is just.; C) Historical factors are secondary to moral imperatives.; D) He has no answer.
\\ \textit{Note:} Nozick argues that as long as each individual's acquisition of holdings is just, the overall distribution can still be considered just, regardless of the relative position of individuals within society, thereby addressing concerns about injustice for those at the bottom.
\item \textbf{Answer:} A
\\ \textit{Options:} A) The Nazi law in question was validly enacted.; B) The court misunderstood the legislation.; C) Fuller misconstrued the purpose of the law.; D) The Nazi rule of recognition was unclear.
\\ \textit{Note:} Hart's response emphasizes that the validity of a law is determined by its adherence to the established legal procedures and rules, regardless of its moral content, thus acknowledging that Nazi law was validly enacted despite its unjust nature.
\item \textbf{Answer:} C
\\ \textit{Options:} A) It is jurisdiction based on the nationality of the offender; B) It is jurisdiction based on where the offence was committed; C) It is jurisdiction based on the nationality of the victims; D) It is jurisdiction based on the country where the legal person was Registered
\\ \textit{Note:} Passive personality jurisdiction refers to the legal principle that allows a state to exercise jurisdiction over a crime committed against its nationals, thus prioritizing the protection of its citizens regardless of where the crime occurred.
\item \textbf{Answer:} D
\\ \textit{Options:} A) Maine; B) Savigny; C) Pound; D) Laski
\\ \textit{Note:} Laski is credited with the phrase "Jurisprudence is the eye of law" because he emphasized the importance of jurisprudence as a critical lens through which legal principles and their application can be understood and interpreted, thus highlighting its foundational role in the study of law.
\end{enumerate}

\section*{True / False}
\begin{enumerate}[label=\arabic*.]
\setcounter{enumi}{5}
\item \textbf{Answer:} False
\\ \textit{Options:} A) True; B) False
\\ \textit{Note:} Jurisdiction is typically territorial, meaning that a nation generally only has legal authority over actions that occur within its own borders, unless specific international agreements or principles apply.
\item \textbf{Answer:} True
\\ \textit{Options:} A) True; B) False
\\ \textit{Note:} H.L.A. Hart is indeed the author of 'The Concept of Law,' a seminal work published in 1961 that profoundly shaped modern legal philosophy and the understanding of legal systems.
\item \textbf{Answer:} False
\\ \textit{Options:} A) True; B) False
\\ \textit{Note:} The statement is false because, according to John Rawls' theory of justice, individuals in the original position would prioritize principles of fairness and equality over absolute liberty to ensure that basic rights and opportunities are available to all members of society.
\item \textbf{Answer:} False
\\ \textit{Options:} A) True; B) False
\\ \textit{Note:} The statement is false because it is H.L.A. Hart, not Holland, who is known for his definition of law as a system of rules applied by the state in the administration of justice.
\item \textbf{Answer:} False
\\ \textit{Options:} A) True; B) False
\\ \textit{Note:} The statement is false because the UN Vienna Declaration of 1993 recognizes civil and political rights as foundational and emphasizes their importance alongside economic, social, and cultural rights, rather than prioritizing third generation rights over them.
\end{enumerate}

\section*{Long Form Response}
\begin{enumerate}[label=\arabic*.]
\setcounter{enumi}{10}
\item \textbf{Answer:} Austin is characterized as a 'naive empiricist' in his approach to law because he emphasizes a pragmatic conception of legal principles over a theoretical framework. He views laws as observable and actionable rules that govern behavior, focusing on their practical application rather than abstract concepts or moral considerations. This perspective aligns with his belief that legal analysis should be grounded in empirical evidence and real-world outcomes, which leads to a more straightforward interpretation of legal norms. By prioritizing the functional aspects of law, Austin's approach tends to overlook the complexities and theoretical underpinnings that often inform legal systems, thereby reinforcing his classification as a naive empiricist.
\\ \textit{Note:} correct because it accurately captures Austin's focus on the practical application of laws as observable rules rooted in empirical evidence, which reflects his characterization as a 'naive empiricist' who prioritizes functionality over theoretical abstraction.
\item \textbf{Answer:} Dworkin criticizes Posner's economic analysis of law by arguing that it erroneously treats wealth as a fundamental value rather than recognizing it as a non-value. He contends that this perspective undermines the moral and ethical dimensions of legal reasoning, particularly in constitutional law, where principles of justice and individual rights should take precedence over mere economic efficiency. By prioritizing wealth maximization, Posner's approach risks neglecting the intrinsic worth of individuals and their rights, which are essential to a just legal system. Dworkin's critique highlights the importance of considering broader values in constitutional interpretation, suggesting that a focus solely on economic outcomes can lead to unjust legal decisions that fail to uphold democratic principles.
\\ \textit{Note:} correct because it accurately captures Dworkin's critique of Posner's economic analysis by highlighting the distinction between wealth as a fundamental value and the moral imperatives that should guide constitutional law, emphasizing the importance of justice and individual rights over economic efficiency.
\end{enumerate}

\vfill
\noindent\textit{End of Answer Key}
\end{document}